%==================== LOI CAM ON ====================
\thispagestyle{empty}
\begin{center}
{\fontsize{17}{22}\selectfont LỜI CẢM ƠN}
\end{center}
\vspace{20pt}

Thời gian làm việc tại KiotViet là quãng thời gian quý báu giúp em tích lũy được nhiều kinh nghiệm thực tế trước khi hoàn thành chương trình đại học. Em xin cảm ơn công ty KiotViet đã tạo điều kiện cho em được tiếp xúc, học hỏi và làm việc trong một môi trường chuyên nghiệp, năng động. Em xin gửi lời cảm ơn chân thành tới anh Nguyễn Mạnh Cường, người trực tiếp hướng dẫn em tại công ty. Anh luôn gợi mở và hướng dẫn em cách tự tìm hiểu bản chất của vấn đề thay vì đưa ra lời giải ngay. Những lần anh yêu cầu em tự giải thích và kiểm tra lại thay đổi trước khi đưa lên hệ thống đã giúp em hình thành thói quen cẩn trọng và có trách nhiệm hơn với công việc của mình.

Em cũng xin cảm ơn các anh chị trong nhóm phát triển sản phẩm khách sạn. Cảm ơn các anh chị đã luôn nhiệt tình hỗ trợ một người mới, kiên nhẫn chỉ ra những điểm chưa tối ưu trong mã nguồn và sẵn sàng chia sẻ kinh nghiệm thực tế giúp em tiến bộ nhanh hơn trong công việc.

Em xin cảm ơn các thầy cô Trường Đại học Công nghệ, Đại học Quốc gia Hà Nội. Những kiến thức nền tảng về cơ sở dữ liệu, cấu trúc dữ liệu và giải thuật cùng công nghệ phần mềm được học tại trường đã giúp em tiếp cận công việc thực tế một cách thuận lợi và vững vàng hơn.

Cuối cùng, em xin gửi lời cảm ơn sâu sắc tới thầy Trần Mạnh Cường vì đã dành thời gian định hướng và hướng dẫn cho em hoàn thành học phần thực tập doanh nghiệp này. Báo cáo khó tránh khỏi những thiếu sót, em rất mong nhận được những ý kiến đóng góp của thầy để báo cáo được hoàn thiện hơn.

Em xin chân thành cảm ơn!

\vspace{24pt}
\begin{flushright}
\begin{minipage}{0.5\textwidth}
\centering
Hà Nội, tháng 9 năm 2026\\
\textit{Sinh viên}\\[26pt]
\textbf{Trần Quang Thành}
\end{minipage}
\end{flushright}

\clearpage
